% GENERATED FILE. Edit canonical lessons, then run npm run generate:books.
% canonical-source-hash: fnv1a64:e73b9076a04d5987
% canonical-lessons: SA-C13-nasika, SA-S212-letter-dha, SA-C13-hridayam

\chapter{The Face, Part Two}
\label{ch:sa-face-nouns-2}
\hlchaptermodality{13}

\begin{chapteropening}
\textbf{By the end of this chapter:} \emph{I can name a nose and a heart in Sanskrit, say which daughter-language nose-word \sk{नासिका} is a relative of rather than a parent of, and close the chapter's twelve nouns against the cousin words this book has already shown.}

Say \sk{हृदयम्}, name the PIE root it shares with heart, cor and kardia, place \sk{हृदयम्} beside the same root's other independent inheritances, and run all twelve of this tranche's nouns back together by group and by gender.
\end{chapteropening}

\section[nāsikā]{\sk{नासिका} (nāsikā) — \textquotedblleft{}nose,\textquotedblright{} kin to two words but parent of neither}

\label{lesson:SA-C13-nasika}



\begin{quote}
Say \textquotedblleft{}mouth.\textquotedblright{} (\emph{Mukham}.) One more part of the face — and a chance to correct a connection this book has not yet made carefully enough.
\end{quote}

\subsection*{You'll want to know: \sk{नासिका}}
\noindent
\begin{tabularx}{\linewidth}{@{}>{\raggedright\arraybackslash}X>{\raggedright\arraybackslash}X>{\raggedright\arraybackslash}X@{}}
\toprule
\textbf{Sanskrit} & \textbf{sound} & \textbf{meaning} \\
\midrule
\sk{नासिका} & \emph{nāsikā} & \textbf{nose} \\
\sk{मम} \sk{नासिका} & \emph{mama nāsikā} & \textbf{my nose} \\
\bottomrule
\end{tabularx}

Feminine, marked by the same \textbf{-\sk{आ}} you will hear again.

\subsection*{The letters in this word}
\textbf{\sk{ना}} is \emph{nā}, \textbf{\sk{न}} under the long-\emph{ā} sign. \textbf{\sk{सि}} is \textbf{\sk{स}} under a short-\emph{i}. \textbf{\sk{का}} closes it, \textbf{\sk{क}} under long-\emph{ā} again: \emph{nā-si-kā}.

\begin{cousinweb}
\textbf{\sk{नासिका}} is built on an older word, \textbf{\sk{नासा}} (\emph{nāsā}), \textquotedblleft{}nose,\textquotedblright{} plus the diminutive suffix \textbf{-\sk{इका}}. \textbf{\sk{नासा}} itself, and the still-older root \textbf{\sk{नस्}} (\emph{nas}) behind it, go back to PIE \textbf{nehas-}, the unbroken root of English \textbf{nose} and Latin \textbf{nasus} ($\to$ \emph{nasal}).

Here is the correction worth making carefully: \textbf{\sk{नासिका}} is \emph{not} the ancestor of Bengali \emph{nāk} or Gujarati \emph{nāk}. Both of those wore straight down from \textbf{\sk{नस्}} itself, through Prakrit \textbf{ṇakka}, by adding a completely different suffix. \sk{नासिका} and \sk{नाक} share a root, the way \textbf{\sk{ग्रह्}} and English \emph{grab} share one — \textbf{relatives, not parent and child}. Reaching for the nearest-looking Sanskrit word is not always reaching for the right ancestor.
\end{cousinweb}

\subsection*{Your turn}
Say these aloud:

\begin{itemize}
\raggedright
  \item \textquotedblleft{}nāsikā\textquotedblright{} — nose
  \item the older word and root beneath it — \emph{nāsā}, \emph{nas}
  \item the word \sk{नासिका} is NOT the parent of — \emph{nāk}
\end{itemize}

\begin{tcolorbox}[breakable,colback=teal!4,colframe=teal!35!black,title={Before you move on}]
What PIE root is \sk{नासिका}'s family from, and name two of its Western cousins. (\textbf{nehas-}; \textbf{nose}, \textbf{nasus}/\emph{nasal}.) Is \sk{नासिका} the ancestor of Bengali \emph{nāk}, or its relative? (\textbf{Relative} — both descend from \sk{नस्}, neither from the other.) What did you learn just before this? (\emph{Mukham}, mouth.)
\end{tcolorbox}
\section[dha]{\sk{ध} — one character, met inside words you already say}

\label{lesson:SA-S212-letter-dha}



\begin{quote}
Before the new one, the ones you have already met: \sk{ह} · \sk{ए} · \sk{ठ}. Say what each of them does.

One character this time. One only — and you have been saying it for pages without knowing which mark on the page it was.
\end{quote}

\subsection*{Script you'll notice: \sk{ध}}
\textbf{\sk{ध}} — \emph{dha}.

It is a \textbf{consonant}, and in this script a consonant is never bare: it comes with an \emph{a} already in it. So it is not \emph{dh}, it is \textbf{dha}.

What it is made of:

\begin{itemize}
\raggedright
  \item an upper spiral ending in a rightward shoulder
  \item a separate lower bowl
  \item a top-to-bottom right stem
  \item the top shirorekhā
\end{itemize}

You already say these, and every one of them has \sk{ध} somewhere inside it:

\begin{itemize}
\raggedright
  \item \textbf{\sk{धन्यवादः}} \emph{dhanyavādaḥ} — thank you
\end{itemize}

\subsection*{Writing: \sk{ध}}
\begin{itemize}
\raggedright
  \item \textbf{1.} start at the upper spiral's inner crossing, curl around the small opening, widen left and down around the outer loop, then sweep right through the shoulder without lifting
  \item \textbf{2.} lift at the left waist and sweep down and around the lower bowl to its right junction
  \item \textbf{3.} lift and draw the right stem top-to-bottom
  \item \textbf{4.} lift and draw the top shirorekhā left-to-right
\end{itemize}

\textbf{Pen lifts: 3.} The pen comes up 3 times and no more.

\begin{quote}
verified four-stroke teaching form; everyday handwriting may join or simplify the loops differently
\end{quote}

\begin{quote}
This is one attested teaching order and not a national standard — handwriting here is taught with school-to-school variation. Source: Opiaterein, ‘Deva-\sk{ध}-order.gif’, strokes 1–4, Wikimedia Commons, 10 May 2009.
\end{quote}

\subsection*{Your turn}
\begin{itemize}
\raggedright
  \item \emph{Look:} at this, and find \sk{ध} in it
\end{itemize}

\begin{quote}
\sk{धन्यवादः}
\end{quote}

\begin{itemize}
\raggedright
  \item \emph{Trace it:} \sk{ध} three times, saying \emph{dha} as you finish each one
  \item \emph{Look:} back at any page of this chapter and find \sk{ध} once more
\end{itemize}

\begin{tcolorbox}[breakable,colback=teal!4,colframe=teal!35!black,title={Before you move on}]
Which character is this — \sk{ध}? What sound does it carry? (\emph{\textbf{dha}}.) Name one word you already say that contains it.
\end{tcolorbox}
\section[hṛdayam]{\sk{हृदयम्} (hṛdayam) — \textquotedblleft{}heart,\textquotedblright{} a root barely disguised anywhere}

\label{lesson:SA-C13-hridayam}



\begin{quote}
Say \textquotedblleft{}nose,\textquotedblright{} \textquotedblleft{}eye,\textquotedblright{} and \textquotedblleft{}mouth.\textquotedblright{} (\emph{Nāsikā}, \emph{akṣi}, \emph{mukham}.) One last word, and it names what this whole chapter has been checking on.
\end{quote}

\subsection*{You'll want to know: \sk{हृदयम्}}
\noindent
\begin{tabularx}{\linewidth}{@{}>{\raggedright\arraybackslash}X>{\raggedright\arraybackslash}X>{\raggedright\arraybackslash}X@{}}
\toprule
\textbf{Sanskrit} & \textbf{sound} & \textbf{meaning} \\
\midrule
\sk{हृदयम्} & \emph{hṛdayam} & \textbf{heart} \\
\sk{मम} \sk{हृदयम्} & \emph{mama hṛdayam} & \textbf{my heart} \\
\bottomrule
\end{tabularx}

Neuter — the family this whole chapter keeps returning to.

\subsection*{The letters in this word}
\textbf{\sk{हृ}} is \emph{hṛ}: \textbf{\sk{ह}} \emph{ha} carrying the vocalic \textbf{\sk{ऋ}}, the same buzzing vowel you first sounded out in \textbf{\sk{गृह्णाति}}. \textbf{\sk{द}} is \emph{da}, \textbf{\sk{य}} is \emph{ya}, \textbf{\sk{म्}} closes it neuter: \emph{hṛ-da-ya-m}.

\begin{cousinweb}
\textbf{\sk{हृदयम्}} is PIE \textbf{kerd-}, and almost nothing about it has changed shape on the way anywhere: English \textbf{heart}, Latin \textbf{cor}, genitive \textbf{cordis} ($\to$ \emph{cordial}, \emph{courage}, \emph{accord}), Greek \textbf{kardia} ($\to$ \emph{cardiac}) — one root, four languages, all still recognisable. The same root anchors Italian \textbf{cuore}, Portuguese \textbf{coração}, and Russian \emph{serdtse}, each an independent inheritance of the identical PIE ancestor rather than a loan from any of the others.

Eastward, Bengali \emph{hridoy} keeps \sk{हृदयम्} whole, a \sk{तत्सम}; its worn-down twin \emph{hiya} survives only in poetry and song, not everyday speech — the rarer pairing, where the carried-whole word is the common one and the worn-down word is the rare one.
\end{cousinweb}

\begin{grammarlens}[title={What you've built: twelve words, three rooms}]
Three groups, twelve nouns: what you would ask for (\emph{pānīyam}, \emph{kṣīram}, \emph{annam}), who you would introduce (\emph{mitram}, \emph{kuṭumbam}, \emph{bhrātā}, \emph{bhaginī}), and your own face (\emph{akṣi}, \emph{karṇaḥ}, \emph{mukham}, \emph{nāsikā}, \emph{hṛdayam}). Almost every one of them ended in the neuter \textbf{-\sk{म्}} you first named beside \textbf{\sk{अन्नम्}} — and almost every one pointed at a word some other language in this project already teaches as its own.
\end{grammarlens}

\subsection*{Your turn}
Say these aloud:

\begin{itemize}
\raggedright
  \item \textquotedblleft{}hṛdayam\textquotedblright{} — heart
  \item four cousins across four languages — heart, cor, kardia, cuore
  \item the Bengali twin kept whole, and the one kept only in poetry — hridoy, hiya
\end{itemize}

\begin{tcolorbox}[breakable,colback=teal!4,colframe=teal!35!black,title={Before you move on}]
Name three languages, besides English, whose word for heart shares \sk{हृदयम्}'s root. (Any three of \textbf{Latin, Greek, Italian, Portuguese, Russian}.) Which Bengali form of \textquotedblleft{}heart\textquotedblright{} is everyday, and which survives only in poetry? (\emph{hridoy}, everyday; \emph{hiya}, poetry.) What gender has run through nearly this entire chapter's nouns? (\textbf{Neuter}.) And what was the very first word of this chapter? (\emph{Pānīyam}, water.)
\end{tcolorbox}
